\documentclass[journal, onecolumn]{IEEEtran}

\usepackage{titlesec}
\usepackage{graphicx}
\usepackage{array}
\usepackage{makecell}
\usepackage{float}
\usepackage{caption}
\usepackage{wrapfig}
\usepackage{url}

\titleformat{\section}
    {\normalfont\bfseries}
    {\thesection}
    {1em}
    {}

\titleformat{\subsection}
    {\normalfont\bfseries}
    {\thesubsection}
    {1em}
    {}

\renewcommand\thesection{\arabic{section}}
\renewcommand\thesubsection{\thesection.\arabic{subsection}}
\renewcommand\thesubsubsection{\thesubsection.\arabic{subsubsection}}

\begin{document}

\title{FireFly Design Document}
\author{Christian Li, Kenny Wang, Kevin Yao, Matin Sobhani, Nick Lafredo, Parsa Rezvani, Ryan Yu}
\date{\today}

\maketitle

\begin{figure*}[!t]
    \centering
    \includegraphics[width=.5\linewidth]{images/UCSC_Logo.png}
    \caption*{}
    \label{fig:UCSCLOGO}
\end{figure*}

\vspace{1in}

\begin{figure*}[!t]
    \centering
    \includegraphics[width=0.25\linewidth]{images/BE_Logo.png}
    \caption*{}
    \label{fig:BELOGO}
\end{figure*}

\clearpage

\tableofcontents

\vspace{1in}

\begin{abstract}
This document presents the design and proof-of-concept for FireFly, a rapidly deployable multi-drone wildfire monitoring system for infrastructure-free terrain. The full design targets autonomous aerial sensing, resilient multi-hop communications, fire localization, and perimeter mapping. The current repository implements the proof-of-concept geolocation pipeline: ESP32 GPS telemetry, a Flask and Leaflet ground station, simulated multi-drone triangulation, YOLOv8 fire detection from video, live MJPEG camera ingest, CAD frame design, bill of materials, and bench-tested ESP-NOW motor and sensor communication.
\end{abstract}

\newpage

% ============================================================
\section*{Executive Summary}
\addcontentsline{toc}{section}{Executive Summary}
% ============================================================

FireFly is designed as a rapidly deployable aerial monitoring system for high wildfire-risk areas where fixed sensors, camera towers, or reliable communications infrastructure are unavailable. The long-term system concept uses a fleet of drones to collect visual, GPS, and environmental observations, relay them through a resilient network, and provide responders with fire alerts, location estimates, and evolving perimeter intelligence.

\vspace{0.5em}

The current implementation is a proof-of-concept for the most important data product: geolocating a fire from drone observations. The repository contains a runnable Python stack that launches a Flask server, displays drones and fires on a Leaflet map, accepts ESP32 GPS telemetry over UDP, simulates a two-drone fleet, and posts fire detections from a YOLOv8 video detector. This demonstrates the ground-station, detection, and triangulation pipeline even though autonomous flight control and full mesh networking are not yet implemented.

\vspace{0.5em}

In the aspirational system, each drone contributes synchronized observations including GPS position, altitude, heading, camera detections, temperature, humidity, smoke or particulate readings, and battery/link health. Multiple drones observing the same ignition point allow the ground station to intersect bearing rays, reject false positives, and estimate a fire location more robustly than a single aircraft.

\vspace{0.5em}

The communication architecture is planned as a split-link system. Low-bandwidth telemetry and alerts use robust packet delivery, while high-bandwidth visual data is prioritized when link quality allows. The current prototype uses Wi-Fi UDP broadcast and optional HTTPS upload for GPS telemetry. MANET, LoRa, and COFDM are design targets for the mature platform rather than demonstrated capabilities in the current repository.

\vspace{0.5em}

Validation is organized in stages. The current evidence includes bench-level ESP-NOW communication and motor/sensor experiments, GPS telemetry firmware, YOLOv8 fire detection scripts, live camera ingest, map visualization, and geometry tests already present in the repo. Future validation will add tethered flight, open-field range testing, calibrated sensor testing, and multi-drone flight trials once the airframe and flight stack are fully integrated.

\clearpage

% ============================================================
\section{Introduction}
% ============================================================

\subsection{Problem Context}
Detecting wildfires in vast, dense forests presents a critical challenge for emergency responders, as early detection is the only way to prevent a small brush fire from becoming an uncontrollable catastrophe. In these environments, a single drone is severely limited by its flight time, battery capacity, and sensor range; it simply cannot scan thousands of acres quickly enough to find a pinpoint heat source before it spreads. Furthermore, a lone drone can only provide a single "line of sight" to a fire, which makes it nearly impossible to determine the exact GPS coordinates or the actual size of the blaze from a distance. Relying on one unit creates a single point of failure where a mechanical issue or a loss of signal in a remote valley results in a total loss of monitoring capability.

To solve this, a mesh network of drones acts as a distributed "eye in the sky," turning a group of individual robots into a cohesive, intelligent sensing platform. By using multiple drones, the search area can be subdivided into a grid, allowing the team to cover a massive forest in a fraction of the time it would take a single unit. When a fire is spotted, the mesh allows multiple drones to converge and view the fire from different angles simultaneously. This multi-perspective data is sent back through the chain of drones—using each other as relays—to the ground station in the city. The central station then uses these intersecting vectors to triangulate the fire's precise location and calculate its growth rate with high accuracy, ensuring that firefighting resources are deployed exactly where they are needed most.

\subsection{Scout Relay Concept}

The primary obstacle in long-distance drone operations is the "line-of-sight" limitation of high-frequency radio signals. In a heavily forested or mountainous environment, physical obstacles like ridges, thick timber, and even the curvature of the Earth act as barriers that swallow data signals. If a single scout drone flies deep into a valley to investigate a heat signature, the terrain will likely sever its direct link to a distant city-based ground station. By strategically positioning relay drones between the scout and the base, the system creates a "daisy chain" that carries the signal over and around these obstacles, ensuring that the scout remains connected no matter how rugged the landscape becomes.

Furthermore, these relay points are essential for maintaining the high bandwidth required for real-time video and complex sensor data. As a radio signal travels, its strength decays rapidly over distance, leading to "noise" and data loss that can cause video feeds to freeze or telemetry to drop. By using a series of relay drones, the data only has to travel a short distance between each "hop." This keeps the signal strength high and the latency low, allowing the artificial intelligence at the ground station to receive a clear, high-definition stream. Without these intermediate points, the data would arrive corrupted or not at all, rendering the scout’s advanced fire-detection sensors useless.

Beyond just signal strength, the relay drones provide a critical layer of operational resilience and "self-healing" for the entire mission. In an active wildfire scenario, conditions change rapidly; smoke plumes can become dense enough to scatter radio waves, or a drone might need to move to avoid extreme heat. A multi-drone mesh allows the network to stay flexible; if one relay point needs to shift its position, the other drones can adjust their spacing to maintain the link. This cooperative behavior ensures a persistent data pipeline, allowing the scouts to focus entirely on the fire while the relays guarantee that the life-saving information they gather actually reaches the decision-makers in the city.

\subsection{MANET}

The planned Mobile Ad-hoc Network (MANET) is the "nervous system" that would transform a group of independent drones into a synchronized sensing fleet. Unlike traditional networks that rely on a fixed router or central tower, a MANET allows each drone to act as both receiver and transmitter. In the unpredictable environment of a forest fire, drones are constantly moving, banking, and changing altitudes to get a better view or to avoid smoke. The MANET layer is intended to manage these shifting positions in real time, automatically calculating the most efficient path for data to travel from the scout to the ground station through available relays.

The most important target characteristic of a MANET in this scenario is self-healing. In a wildfire mission, the risk of losing a drone to technical failure, low battery, or extreme environmental conditions is high. In a rigid network, the loss of a single relay drone can break the communication chain. A decentralized mesh should let the remaining drones reroute data through other nearby units as long as a physical path to the base remains.

Furthermore, MANET technology allows for dynamic scalability, which is essential when the scale of a fire grows unexpectedly. If a fire spreads and more drones are deployed from the city to assist, they can join the MANET instantaneously without needing a manual configuration or a reboot of the system. The network simply expands to include the new "nodes," automatically incorporating them into the relay chain or the search grid. This "plug-and-play" nature of ad-hoc networking allows emergency responders to scale their surveillance efforts in real-time as the situation on the ground evolves.

Finally, using a MANET is important because it reduces the "overhead" and interference that usually plagues large-scale wireless systems. By using intelligent routing protocols, the drones can "gossip" with their immediate neighbors to determine which path is the least congested or has the strongest signal. This prevents the airwaves from becoming cluttered with redundant transmissions, ensuring that the high-priority triangulation data and AI video feeds reach the ground station with minimal latency. In a life-or-death firefighting scenario, the milliseconds saved by an optimized MANET path can be the difference between catching a fire early and losing control of the forest.

\subsection{LoRa and COFDM}

LoRa, short for Long Range, is a foundational component of our communication strategy, utilizing a spread-spectrum modulation technique that allows for connectivity over extreme distances. Its primary strength lies in its "receive sensitivity," meaning it can maintain a link even when the signal is incredibly faint or obstructed by dense forest canopies. For emergency operations, this provides a nearly unbreakable "lifeline" for the most critical data points: the drone’s GPS coordinates, battery health, and high-priority fire detection alerts.

However, the physical laws of radio frequency dictate a strict trade-off between range and speed. Because LoRa is optimized for distance and low power consumption, it has a very narrow data bandwidth—equivalent to sending simple text strings rather than complex media. While it is the ultimate tool for ensuring we never lose track of a drone, it is fundamentally incapable of transmitting the high-definition video streams required for modern AI-driven analysis. If used in isolation, the system would be "aware" of a drone's location but "blind" to the visual details of the fire.

This limitation creates a significant hurdle for remote sensing: a mission that only provides coordinates without visual confirmation lacks the context needed for high-level decision-making. Investors must recognize that relying on a single radio protocol in a vast forest creates a bottleneck that renders advanced AI software useless. To bridge the gap between "long range" and "high quality," a more sophisticated, multi-layered approach is required to meet the demands of professional firefighting agencies.

This motivates the planned Split-Link Architecture. This design philosophy separates communication into two distinct "pipes" operating simultaneously. Rather than forcing a single radio to compromise on both speed and distance, the mature system would use specialized hardware for low-bandwidth telemetry and higher-bandwidth video. This dual-path strategy keeps the "nervous system" of the fleet--telemetry, GPS, bearing, and alert data--independent from the "eyes" of the fleet--the video feed.

For the future high-bandwidth path, the system evaluates COFDM (Coded Orthogonal Frequency Division Multiplexing) because OFDM-style links can tolerate multipath better than simple single-carrier video links. In a forest, radio signals are reflected off trees and terrain, which can fragment video. A COFDM video path is therefore treated as a planned hardware upgrade for non-line-of-sight video, not as a capability demonstrated by the current repository.

In the current prototype, the video path is implemented with recorded fire videos, ESP32-CAM or MJPEG stream URLs, OpenCV frame capture, and YOLOv8 inference running on the ground-station computer. The proof-of-concept therefore validates detection and map posting, while future RF hardware must still be selected and tested for real outdoor range, bandwidth, latency, and multipath performance.

The Split-Link Architecture is intended to solve range and video quality as separate engineering problems. If the high-bandwidth video link degrades, the lower-bandwidth telemetry path should continue to carry GPS, heading, battery state, and fire-alert packets. This supports a "fail-functional" mission mode in which the operator may lose video detail but still retains the data required to track drone position and estimate fire location.

Ultimately, dual-protocol hardware is a future design direction for operational resilience. The implemented repository currently uses Wi-Fi UDP telemetry and HTTP APIs, while the MANET, LoRa, and COFDM portions remain requirements for the next hardware iteration.

\subsection{Automation Concept}
To maximize operational efficiency and ensure rapid response, the full FireFly concept uses a dynamic automation workflow that transitions through three distinct phases: the Lawnmower Search, the Relay Chain, and the Orbital Ring. This tiered approach is a design target for future autonomous flight software. The current repository simulates drone movement and validates how the ground station reacts to detection and telemetry data.

In the initial "Search Phase," the fleet executes a Lawnmower pattern, which is the industry standard for 100 percent visual coverage. The mission area is divided into precise parallel "lanes," and each drone is autonomously assigned a specific path. This systematic grid ensures that we are not just flying aimlessly, but are instead "scanning" the forest floor with mathematical certainty. By subdividing the area among multiple drones, we can search thousands of acres in minutes—a task that would take hours for a single aircraft or ground crew.

As the search expands farther from the ground station, the future system would transition into the Relay Chain configuration. Not all drones are sent to the front line; instead, automation software would position selected units as high-altitude signal repeaters. This creates a resilient data bridge for future video and telemetry links. This self-positioning relay logic is not currently implemented in the repository.

The moment a scout drone’s onboard AI detects a fire, the fleet enters the Orbital Ring phase. Rather than staying in a fixed grid, nearby drones are "called" to the detection site to surround the fire from multiple cardinal directions. This orbital formation is crucial for the triangulation math performed at the ground station. By viewing the heat source from three or four different angles simultaneously, the system can intersect those vectors to pinpoint the fire’s GPS coordinates and estimate its total acreage with far greater accuracy than a single observation ever could.

For the customer, this three-stage automation means the entire mission—from search to discovery to precise location tracking—is handled with minimal human intervention. The drones "decide" how to move based on the needs of the network and the urgency of the fire. This workflow provides a professional-grade solution that combines the speed of a broad search with the tactical precision of a localized orbital scan, ensuring that firefighting teams have the exact data they need the moment a fire begins to spread.


\vspace{1em}

Wildfires exert a devastating, multi-layered impact on communities, beginning with the immediate destruction of physical infrastructure, homes, and essential utilities that can displace thousands and cripple local economies. Beyond the visible charred landscape, these events trigger severe public health crises, as fine particulate matter from smoke travels hundreds of miles, exacerbating respiratory and cardiovascular conditions for both local residents and distant populations. The long-term recovery is often hindered by "burn scars" on the environment, where the loss of vegetation and changes in soil chemistry significantly increase the risk of flash flooding and debris flows during subsequent rainstorms. Furthermore, the psychological toll—manifesting as trauma, chronic stress, and a lost sense of security—can linger for years, fundamentally altering the social fabric and demographic makeup of a community as residents struggle with the financial and emotional burdens of rebuilding or relocating.

\vspace{1em}

The best way to stop these devastating effects is to act on the problem before the fire becomes uncontainable. Early wildfire detection is the most critical factor in transforming a potential catastrophe into a manageable incident, as it allows fire suppression teams to intervene during the "initial attack" phase when the fire’s perimeter is still small and its energy output is low. By identifying an ignition point within minutes, responders can contain flames before they reach the crowns of trees or transition into high-intensity firestorms that create their own weather patterns. This rapid response time not only drastically reduces the acreage burned and the associated carbon emissions but also prevents the fire from reaching the "Wildland-Urban Interface" (WUI), where the proximity to human structures makes containment significantly more dangerous and costly. Ultimately, catching a fire in its infancy preserves critical infrastructure, maintains local air quality, and saves lives by providing residents with more adequate lead time for evacuations.

\vspace{1em}

A key limitation in existing wildfire monitoring systems is the lack of persistent, high resolution observation during the earliest stages of ignition, particularly in remote or infrastructure limited environments. Satellite based systems provide broad coverage but suffer from temporal latency, while ground based sensors are sparse and often unavailable in high risk wilderness regions. As a result, early fire growth and boundary formation often occur without sufficient real time spatial data.

\vspace{1em}

FireFly is our mobile drone early fire detection system. It is designed to support early ignition detection and real time wildfire intelligence for first responders. The platform operates as a multi-drone autonomous sensing network that can be deployed to a geospatial region of interest using approximate coordinates. Once deployed, each drone uses live visual input combined with onboard inference to navigate toward active fire zones for closer inspection and continuous monitoring.

\vspace{1em}

FireFly addresses the shortcomings of fixed monitoring infrastructure by enabling mobile aerial tracking of early stage fire events. In the intended operational configuration, a fleet of three or more drones collaborates to localize fire sources through multi-point observation and triangulation. Each drone contributes synchronized visual data and localized environmental measurements. The current prototype demonstrates this concept with two simulated drones, ESP32 GPS telemetry, a server-side triangulation module, and a map dashboard that plots fire estimates.

\vspace{1em}

The system prioritizes adaptive fire boundary mapping as its core functional output. As the fire evolves, FireFly continuously updates perimeter estimates and directional spread models, enabling the identification of high risk zones, propagation direction, and relative intensity changes over time. This persistent mapping capability provides responders with continuously updated situational awareness rather than static snapshots, supporting more informed decision making during rapidly changing fire conditions.

\subsection{Need Statement}
Remote and high fire-risk areas lack the sensor infrastructure necessary to monitor environmental conditions during the critical pre-fire window, leaving fire response teams without the real-time data needed to act early.

\subsection{Goal Statement}
Provide fire response teams with real-time temperature, humidity, and smoke readings from high-risk, infrastructure-free terrain by deploying a rapidly mobilized monitoring system that transmits real-time data to a ground station and alerts operators when fire risk thresholds are crossed.

\subsection{Personas}
Individuals with experience levels ranging from novice to expert in fire monitoring and environmental fieldwork.

\subsubsection*{End Users}

\begin{wrapfigure}{l}{0.2\textwidth}
    \includegraphics[width=0.15\textwidth]{images/marcus.png}
\end{wrapfigure}
\textbf{Marcus} \\
\underline{\textit{Role:}} Wildfire Operations Coordinator (professional first responder) \\
\underline{\textit{Background / Context:}} Works with a state forestry division managing prescribed burns and post-fire monitoring. Needs something fast that he can hand off to field crews without extensive training. \\
\underline{\textit{Goals:}} Rapid deployment of sensor-equipped drones into areas too dangerous or inaccessible for personnel during the high-risk pre-fire window. Wants real-time smoke and temperature data without putting people in harm's way. \\
\underline{\textit{Behaviors:}} Works under time pressure, values reliability over features, and coordinates with multiple agencies. Skeptical of unproven technology. Practical: if it breaks in the field, it needs to be fixable on-site with minimal tools. \\
\underline{\textit{Expectations:}} A rugged, simple-to-operate platform that can be deployed at a moment's notice. Wants clear threshold alerts he can act on immediately.

\vspace{1em}

\begin{wrapfigure}{l}{0.2\textwidth}
    \includegraphics[width=0.18\textwidth]{images/elena.png}
\end{wrapfigure}
\textbf{Elena} \\
\underline{\textit{Role:}} Ecology Graduate Researcher (academic user) \\
\underline{\textit{Background / Context:}} Studying fire risk and ecosystem conditions at a land-grant university. Needs to collect environmental sensor data across wilderness areas where no permanent monitoring infrastructure exists. Currently hikes in with handheld sensors, which limits how often and how broadly she can sample. \\
\underline{\textit{Goals:}} Automate pre-fire data collection over defined survey areas. Wants live temperature, humidity, and smoke readings logged without requiring physical presence at each site. \\
\underline{\textit{Behaviors:}} Comfortable with technical documentation, participates in academic drone and environmental monitoring communities, and documents everything for publication. Willing to configure and customize but does not have time to build from scratch. \\
\underline{\textit{Expectations:}} A reliable platform with a clear data pipeline she can connect to her existing research database, and results trustworthy enough to cite.

\vspace{1em}

\subsubsection*{Experts / Consultants}

\begin{wrapfigure}{l}{0.2\textwidth}
    \includegraphics[width=0.18\textwidth]{images/anita.png}
\end{wrapfigure}
\textbf{Dr. Anita Rowe} \\
\underline{\textit{Role:}} Environmental Sensor Systems Engineer \\
\underline{\textit{Background / Context:}} Works at a national lab developing atmospheric sensor arrays for wildfire monitoring. Has deep expertise in what data is actually actionable in pre-fire environments and where sensors fail under heat, humidity, and particulate exposure. \\
\underline{\textit{Goals:}} Validate that FireFly's sensor suite captures meaningful data for fire risk assessment and identify unreliable component choices before they are locked into the design. \\
\underline{\textit{Behaviors:}} Data-driven and skeptical of consumer-grade sensors in harsh field conditions. Will probe calibration drift, operating temperature limits, and whether threshold alert logic reflects real-world fire risk indicators. \\
\underline{\textit{Expectations:}} A platform sound enough in its sensor selection and alert logic that field results can be trusted and the methodology can be defended to other scientists.

\vspace{1em}

\subsubsection*{Client}

\begin{wrapfigure}{l}{0.2\textwidth}
    \includegraphics[width=0.18\textwidth]{images/usfs.png}
\end{wrapfigure}
\textbf{U.S. Forest Service Regional Office} \\
\underline{\textit{Role:}} Federal Land and Fire Management Agency \\
\underline{\textit{Background / Context:}} Responsible for fire prevention, monitoring, and response across millions of acres of national forest, much of it roadless wilderness with no existing sensor infrastructure. Currently relies on satellite imagery, fire lookout towers, and aerial patrols, all of which have coverage gaps and response delays during the highest-risk windows. \\
\underline{\textit{Goals:}} Evaluate FireFly as a rapidly deployable supplement to existing monitoring systems for areas where infrastructure cannot operate. Interested in whether the platform can reliably trigger alerts before a fire establishes. \\
\underline{\textit{Behaviors:}} Evaluates tools based on operational reliability and ease of use for non-specialist field crews. Requires demonstrated performance data before recommending adoption. \\
\underline{\textit{Expectations:}} A working prototype that proves the core concept: zero-infrastructure deployment, live sensor transmission, and threshold-based alerting, with documentation clear enough for field personnel to operate without engineering support.

\subsection{Existing Research}
Current wildfire detection relies on satellite imagery, fixed camera networks, manned aircraft, and ground reports. NASA FIRMS distributes MODIS and VIIRS active-fire products; MODIS fire pixels are approximately 1 km at nadir, while the VIIRS I-band product provides nominal 375 m active-fire pixels and improved mapping of smaller fires. These systems are valuable for regional awareness, but satellite overpass timing, pixel size, clouds, smoke, and canopy obstruction can limit early ignition confirmation. Fixed camera networks such as ALERTCalifornia provide real-time visual monitoring at large scale, but each camera remains constrained by tower placement, terrain, weather, and line of sight.

\vspace{0.5em}

FireFly addresses this gap by providing a self-contained, rapidly deployable aerial sensing platform. The full system is intended to supplement satellites and fixed towers with lower-altitude, taskable observations in infrastructure-free terrain. The current prototype demonstrates the software pipeline needed to turn drone telemetry and visual detections into map-based fire coordinates.

\subsection{Sustainability Statement}

The FireFly project is committed to the ``Design for Environment'' philosophy. Our goal is to provide a high-tech solution for wildfire mitigation while ensuring that the lifecycle of the drone itself does not offset its environmental benefits.

\vspace{1em}

\textbf{Environmental Stewardship:}
The primary purpose of FireFly is the preservation of biodiversity and carbon sinks by preventing catastrophic wildfires. By providing early-warning data in the pre-fire window, we reduce the need for large-scale, carbon-intensive emergency responses such as heavy tanker aircraft and ground fleets. Furthermore, by utilizing electric propulsion over internal combustion, we ensure zero-emission operation during monitoring flights.

\vspace{1em}

\textbf{Material Selection and Circularity:}
To minimize our manufacturing footprint, we have opted for 3D-printed components using \textbf{PLA}. PLA is a bio-based, industrially compostable polymer derived from renewable resources like corn starch. By using additive manufacturing, we significantly reduce material waste compared to traditional subtractive machining.

\vspace{1em}

\textbf{Life Cycle and Battery Management:}
We recognize that Lithium-Polymer (LiPo) batteries pose a chemical risk if disposed of improperly. Our operational plan includes a strict end-of-life protocol where degraded cells are discharged and taken to certified e-waste recycling centers. We have designed the electronics with standardized connectors to ensure that sensors and flight controllers can be harvested and reused in future iterations of the platform, extending the functional life of the high-value silicon components.

\vspace{1em}

\textbf{Social and Safety Impact:}
Sustainability includes the protection of human capital. FireFly enhances the safety of first responders by providing situational awareness in hazardous terrains without risking human life. By creating a rapidly deployable tool, we extend access to advanced fire-monitoring capability to smaller forestry divisions and academic researchers who operate outside the reach of satellite-based or manned aerial surveillance programs.

\vspace{1em}

\textbf{Ethics Statement:}
As engineers, we recognize that our work has a direct impact on public safety and the environment. We are committed to upholding the highest standards of honesty, transparency, and responsibility throughout the design and development of FireFly. Our system is intended solely to support fire prevention and responder safety. We will not misrepresent the capabilities of the prototype, will clearly communicate the limitations of sensor accuracy and flight duration, and will ensure that any data collected during testing is handled responsibly. We adhere to the NSPE Code of Ethics, holding paramount the safety, health, and welfare of the public in all design decisions.

\subsection{Operational Scenarios for Wildfire Monitoring}
FireFly is designed to perform in three main scenarios to assist first responders quickly and ecological researchers in monitoring and analyzing active wildfire threats.

\vspace{1em}

\textbf{Active Fire Detection:}
FireFly is intended to serve as a mobile monitoring node in blind spot regions where satellites or fixed towers are delayed or blocked by topography. The current prototype uses YOLOv8 image recognition on recorded video or live MJPEG streams to detect fire and smoke, then posts detections to the map for operator review.

\vspace{1em}

\textbf{Fire Propagation and Heading Analysis:}
In the full system, repeated geolocated observations, calibrated camera data, and environmental gradients would support estimates of direction, speed, and intensity change. The repository currently supports heading-based fire localization and stores fire markers over time; propagation modeling remains future work.

\vspace{1em}

\textbf{Perimeter Mapping and Size Assessment:}
In the full system, multiple drones can be deployed around the boundary of a detected fire to generate a geofenced perimeter and area estimate. The current proof-of-concept estimates fire points from bearing intersections and bounding-box size approximations; accurate perimeter mapping requires additional flight, camera, and calibration validation.

% ============================================================
\section{Design}
% ============================================================

\subsection{Design Objectives}
The objective of FireFly is to develop a robust aerial monitoring platform capable of rapid deployment in remote, high fire-risk environments. The full system is designed around five engineering goals: maintain useful coverage without fixed infrastructure, detect fire and smoke signatures from aerial imagery, localize ignition points using multiple drone observations, transmit low-latency telemetry to a ground station, and give operators a clear map-based view of drones, detections, and fire estimates. The current repository proves the geolocation and ground-station portions of this design, while full autonomous flight, outdoor relay networking, calibrated environmental alerting, and live multi-drone perimeter tracking remain next-iteration work.

\subsection{System Requirements and Acceptance Criteria}

\begin{table}[H]
    \centering
    \small
    \begin{tabular}{|p{0.12\textwidth}|p{0.26\textwidth}|p{0.18\textwidth}|p{0.20\textwidth}|p{0.10\textwidth}|}
        \hline
        \textbf{ID} & \textbf{Requirement} & \textbf{Target} & \textbf{Verification Method} & \textbf{Status} \\
        \hline
        REQ-01 & Mission endurance / usable flight time & 5--8 min for Option A hardware estimate & eCalc estimate, then tethered and field flight timing & Estimated \\
        \hline
        REQ-02 & Telemetry range & Local Wi-Fi for prototype; multi-hop relay for final system & UDP packet range test and mesh relay field test & Prototype only \\
        \hline
        REQ-03 & Telemetry update interval & 1 Hz GPS in fleet simulator; firmware sends every valid GPS fix & Server logs, sequence numbers, and packet-drop script & Implemented \\
        \hline
        REQ-04 & Sensor readout accuracy / stability & Temperature and humidity values readable over I2C; calibration required before field use & SHTC3 serial output and reference-instrument comparison & Bench tested \\
        \hline
        REQ-05 & Alert generation delay from visual fire detection & 3 s sustained fire in recorded-video detector; 2 s sustained fire in MJPEG ingest & Controlled video or injected stream detection & Implemented in software \\
        \hline
        REQ-06 & Deployment and setup time & One-command local demo after dependencies are installed & Run launcher with server, simulator, and optional camera discovery & Implemented for demo \\
        \hline
        REQ-07 & Maximum takeoff mass / payload margin & Option A baseline approximately 190--230 g, payload comfort 0--75 g & Mass scale and thrust model comparison & Estimated \\
        \hline
    \end{tabular}
    \caption{System requirements and acceptance criteria.}
    \label{tab:requirements_template}
\end{table}

\subsection{System Architecture Overview}

FireFly is organized as a sensing fleet, communications layer, and ground-station intelligence layer. In the full system, each drone carries flight-control hardware, GPS, camera, environmental sensors, and communications hardware. The ground station receives telemetry and detections, fuses multi-drone observations, and presents fire intelligence to the operator. The current repository implements a proof-of-concept version of this architecture using ESP32 GPS telemetry firmware, Python simulation, Flask APIs, Leaflet map visualization, OpenCV video ingest, and YOLOv8 fire detection.

\subsubsection*{Subsystem Inventory}
\begin{itemize}
    \item Airframe / mechanical structure: 3D-printed X-frame concept with v5 CAD iteration, motor pods, battery strap slots, and serviceable central bay.
    \item Propulsion stack: SunnySky R1106-class motors, tri-blade propellers, 3S LiPo battery, and 4-in-1 ESC selected for the intended flight build; bench testing used an H-bridge and brushed DC motor while waiting for final propulsion hardware.
    \item Flight stabilization stack: planned flight-controller integration; not demonstrated by the current geolocation-focused repository.
    \item Environmental sensing stack: SHTC3 temperature/humidity firmware path is implemented; smoke/gas, CO2, UV/IR, and particulate sensing are design targets requiring final sensor selection and calibration.
    \item Wireless telemetry link: current GPS telemetry uses Wi-Fi UDP broadcast on port 4210 and optional HTTPS upload; ESP-NOW is bench-tested for controller-to-drone communication; ESP-MESH, LoRa, and COFDM are future architecture elements.
    \item Ground station: Flask backend, Leaflet map, JavaScript dashboard, Supabase-compatible drone registration path, camera URL registration, YOLOv8 video detector, MJPEG ingest, and triangulation module.
\end{itemize}

\subsubsection*{End-to-End Data Path}
\begin{table}[H]
    \centering
    \small
    \begin{tabular}{|p{0.18\textwidth}|p{0.18\textwidth}|p{0.18\textwidth}|p{0.16\textwidth}|p{0.18\textwidth}|}
        \hline
        \textbf{Stage} & \textbf{Source} & \textbf{Destination} & \textbf{Format / Protocol} & \textbf{Notes} \\
        \hline
        1 & GPS receiver & ESP32 GPS firmware & UART, NMEA GGA & Parses latitude, longitude, altitude, satellites, and HDOP \\
        \hline
        2 & ESP32 firmware & Ground station LAN & UDP JSON broadcast, port 4210 & Fields include id, seq, lat, lon, alt\_m, heading\_deg, fire\_detected, sats, hdop, ts\_ms \\
        \hline
        3 & Flask GPS listener & In-memory drone state and optional database & Python dictionary / SQLAlchemy models & Drops stale drones after 30 s \\
        \hline
        4 & Flask server & Browser dashboard & HTTP JSON via \verb|/api/state| & JavaScript polls every 500 ms and updates the map \\
        \hline
        5 & YOLO detector or triangulator & Ground station fire list & HTTP JSON via \verb|/api/fire| & Fire marker includes coordinates, confidence, area, and source when available \\
        \hline
    \end{tabular}
    \caption{Current proof-of-concept data path.}
    \label{tab:architecture_template}
\end{table}

\subsection{Mass Budget}

\begin{table}[H]
    \centering
    \small
    \begin{tabular}{|p{0.28\textwidth}|p{0.12\textwidth}|p{0.14\textwidth}|p{0.14\textwidth}|p{0.18\textwidth}|}
        \hline
        \textbf{Component} & \textbf{Qty.} & \textbf{Mass Each} & \textbf{Total Mass} & \textbf{Measured / Estimated} \\
        \hline
        Frame & 1 & 30 g & 30 g & Estimated from design notes \\
        \hline
        SunnySky R1106 motors & 4 & 7.5 g & 30 g & Vendor/design estimate \\
        \hline
        Propeller set & 1 & 5 g & 5 g & Estimated from design notes \\
        \hline
        Battery, 850 mAh 3S LiPo & 1 & 69 g & 69 g & Estimated from design notes \\
        \hline
        ESP32 controller & 1 & 9 g & 9 g & Estimated from design notes \\
        \hline
        ESC, wiring, sensors, mounts & 1 lot & Not isolated & 47--87 g & Budgeted to reach 190--230 g AUW \\
        \hline
        Optional payload margin & 1 lot & 0--75 g & 0--75 g & Option A payload comfort estimate \\
        \hline
        Total all-up weight & & & 190--230 g & Estimated; needs scale measurement \\
        \hline
    \end{tabular}
    \caption{Prototype mass budget from repo design notes.}
    \label{tab:mass_budget_template}
\end{table}

\subsection{Sensor and Alert Logic}

\subsubsection*{Sensor Baseline}
\begin{itemize}
    \item Current implemented sensors/data sources: GPS over UART, SHTC3 temperature/humidity over I2C in the motor receiver prototype, potentiometer command input for motor bench testing, recorded or live MJPEG video for YOLO fire detection.
    \item Current software-only inputs: simulated drone positions, headings, and fire visibility from \verb|simulate_fleet.py|.
    \item Future-work sensors: smoke/gas detector, CO2 sensor, particulate sensor, UV/IR flame sensor, LiDAR/range sensor, battery monitor, and calibrated thermal imaging.
    \item Sampling/update rates: simulated fleet broadcasts at 1 Hz; GPS firmware broadcasts each valid GGA fix; motor command prototype sends every 50 ms; SHTC3 telemetry is sent every 2 s; dashboard polling is 500 ms.
    \item Expected operating range: environmental sensor limits and calibration targets must be finalized from selected datasheets before field deployment.
\end{itemize}

\subsubsection*{Alert Thresholds and Current Logic}
\begin{table}[H]
    \centering
    \small
    \begin{tabular}{|p{0.20\textwidth}|p{0.20\textwidth}|p{0.16\textwidth}|p{0.18\textwidth}|p{0.16\textwidth}|}
        \hline
        \textbf{Signal} & \textbf{Condition} & \textbf{Threshold} & \textbf{Filtering / Time Window} & \textbf{Alert Output} \\
        \hline
        Temperature & Absolute temperature or rapid rise & To be calibrated for final sensor & Moving average and rate-of-change planned & Risk score and operator warning \\
        \hline
        Humidity & Low relative humidity combined with heat & To be calibrated for final sensor & Combined with temperature and smoke & Risk score modifier \\
        \hline
        Smoke / particulate / gas & Elevated smoke or gas reading & Sensor-specific threshold required & Debounce window planned & Fire-risk alert \\
        \hline
        YOLO fire class & Fire object detected in video & Confidence 0.4 in recorded-video script & Fire must persist for 3 s & POST to \verb|/api/fire| \\
        \hline
        MJPEG fire/smoke class & Fire or smoke in live stream & Confidence 0.25 in ingest worker & Fire/smoke must persist for 2 s; 5 s cooldown & POST to \verb|/api/fire| \\
        \hline
    \end{tabular}
    \caption{Alert logic for current software and planned environmental sensing.}
    \label{tab:alert_threshold_template}
\end{table}

\subsubsection*{Calibration and Environmental Limits}
\begin{itemize}
    \item Calibration reference used: no completed calibration log is present in the repo; final temperature, humidity, smoke, and gas thresholds must be validated against reference instruments.
    \item Pre-flight sensor validation procedure: confirm GPS fix, verify nonzero telemetry updates, check camera stream availability, verify environmental sensor readings against ambient conditions, and confirm dashboard receipt before launch.
    \item Known limitations: GPS jitter affects heading when nearly stationary, smoke may dilute in wind, heat soak may bias onboard temperature sensors, and video detections depend on lighting, viewing angle, model confidence, and camera calibration.
    \item False-positive and false-negative mitigation: require sustained visual detections, combine environmental signals with image confidence, use multi-drone bearing agreement, and allow operator confirmation through the map and camera feed.
\end{itemize}

\subsection{Safety, Reliability, and Operational Constraints}

\subsubsection*{Risk Register}
\begin{table}[H]
    \centering
    \small
    \begin{tabular}{|p{0.22\textwidth}|p{0.18\textwidth}|p{0.16\textwidth}|p{0.18\textwidth}|p{0.14\textwidth}|}
        \hline
        \textbf{Failure / Hazard} & \textbf{Cause} & \textbf{Impact} & \textbf{Mitigation / Failsafe} & \textbf{Residual Risk} \\
        \hline
        Link loss & Wi-Fi range, terrain, packet loss, stale stream & Ground station loses drone position or camera feed & Dashboard drops stale drones after 30 s; future mesh/relay link required & Medium \\
        \hline
        Low battery / brownout & LiPo discharge under propulsion load & Loss of control or telemetry reset & Future battery monitor, return-to-home threshold, and preflight voltage check & High until flight stack is integrated \\
        \hline
        Sensor fault / drift & GPS jitter, I2C read failure, heat soak, camera calibration error & Wrong heading, poor fire localization, false alert & CRC checks where available, calibration, sustained detection windows, multi-drone confirmation & Medium \\
        \hline
        Propulsion failure & Motor, ESC, wiring, propeller, or frame damage & Crash or uncontrolled landing & Bench motor tests, tethered flight before open field, inspection checklist & High until flight testing is complete \\
        \hline
        Structural damage / crash & Hard landing, arm fracture, motor pod mismatch & Loss of airframe or payload & One-piece reprintable frame, visual inspection, conservative test envelope & Medium \\
        \hline
        Unauthorized command injection & Unauthenticated ESP-NOW command packets & Motor command spoofing in bench prototype & Validate peer MAC, clamp PWM values, enable encryption or authenticated protocol before flight & High until remediated \\
        \hline
    \end{tabular}
    \caption{Risk and reliability register.}
    \label{tab:risk_template}
\end{table}

\subsubsection*{Operational Constraints}
\begin{itemize}
    \item Environmental limits: prototype testing should use calm wind, clear visibility, open terrain, and conservative altitude until flight stability and telemetry range are measured.
    \item Regulatory constraints: outdoor flight must comply with FAA small UAS rules, including visual line of sight, operating altitude, airspace authorization, and pilot responsibilities.
    \item Field safety procedures: use a defined exclusion radius, propeller-off bench tests, tethered first flights, LiPo-safe charging and storage, preflight communications checks, and documented recovery procedures.
\end{itemize}

\newpage

\subsection{Aesthetic Prototype}

The FireFly aesthetic prototype follows a field-ready design language: a compact enclosed central body with rounded corners, symmetric X-arm geometry, and cylindrical motor pods at each corner. The overall form emphasizes structural clarity and quick visual inspection rather than cosmetic styling. Rounded exterior edges reduce snag points during transport, while the open-top bay and side windows communicate where electronics, wiring, and battery hardware are intended to be serviced.

\begin{figure}[H]
    \centering
    \includegraphics[width=0.46\textwidth]{images/drone_design_5_top.png}
    \hfill
    \includegraphics[width=0.46\textwidth]{images/drone_design_5_corner.png}
    \caption{FireFly CAD aesthetic prototype, iteration v5 (top and isometric views).}
    \label{fig:drone_v5_aesthetic}
\end{figure}

The prototype evolved through multiple print-test iterations. Early versions validated basic packaging and motor placement, while later versions introduced a cleaner body-to-arm transition, larger internal volume, and improved vent/window geometry for lighter weight and faster access to internals. This iterative process preserved a consistent visual identity while converging on a more manufacturable and maintainable frame.

\begin{figure}[H]
    \centering
    \includegraphics[width=0.85\textwidth]{images/drone_design_5_compare_1.jpeg}
    \caption{Physical prototype progression (v1--v3) that informed the final v5 frame direction.}
    \label{fig:drone_proto_progression}
\end{figure}

\subsection{Design for Manufacture, Assembly, Maintenance}

\textbf{Design for Manufacture (DFM):}
The current frame is designed as a single-piece 3D print to minimize part count and remove alignment errors. The v5 CAD envelope is approximately $100 \times 72 \times 24$ mm with $2.5$ mm floor and wall thicknesses, balancing print speed, stiffness, and mass. Arms are sized at $10.0$ mm width and $5.4$ mm thickness to maintain rigidity under motor thrust and landing loads. Motor pods are dimensioned for SunnySky R1106-class motors with built-in tolerance ($19$ mm outer diameter, $15.0$ mm inner diameter, $17$ mm height, $2.2$ mm mounting holes on a $9$ mm bolt-circle).

\vspace{0.5em}

\textbf{Design for Assembly (DFA):}
Assembly features are integrated directly into the print geometry. The frame includes four battery strap slots in the floor and continuous wire-routing paths from motor pods into the electronics bay. This allows a straightforward assembly order: install motors in pods, route motor leads through arm channels, mount the controller board, then secure the battery and sensor harness. Because most interfaces are geometric rather than adhesive, rework is possible without damaging the frame.

\vspace{0.5em}

\textbf{Design for Maintenance (DFMn):}
Maintenance access is provided by the open top cavity, side-wall windows, and pod-to-arm wire notches. These features allow quick inspection of solder joints, connector seating, and wire wear without full teardown. Electronics can be removed independently from propulsion hardware, and battery straps can be replaced in the field with standard hook-and-loop straps. In case of structural damage, the one-piece architecture supports rapid replacement by reprinting the frame using the same CAD file, preserving dimensional repeatability across repairs.

\subsection{Block Diagrams}

\begin{figure}[H]
    \centering
    \includegraphics[width=1.0\textwidth]{images/Final_Schematic.png}
    \caption{A high-level schematic of the FireFly system.}
    \label{fig:schematic_design}
\end{figure}

\subsection{Wiring Diagrams}

\begin{figure}[H]
    \centering
    \includegraphics[width=0.75\textwidth]{images/Circuit Diagram.png}
    \caption{Circuit Diagram}
    \label{fig:circuit_diagram}
\end{figure}

\subsection{State Transition Diagrams}

\begin{figure}[H]
    \centering
    \includegraphics[width=0.75\textwidth]{images/State Transition Diagram.png}
    \caption{State Transition Diagram}
    \label{fig:state_transition}
\end{figure}

\subsection{Interface Definitions}

\begin{table}[H]
    \centering
    \scriptsize
    \begin{tabular}{|p{0.16\textwidth}|p{0.14\textwidth}|p{0.14\textwidth}|p{0.14\textwidth}|p{0.12\textwidth}|p{0.12\textwidth}|p{0.12\textwidth}|}
        \hline
        \textbf{Interface} & \textbf{Source} & \textbf{Destination} & \textbf{Physical Layer} & \textbf{Voltage Level} & \textbf{Update Rate} & \textbf{Notes} \\
        \hline
        GPS UART & GPS module & ESP32 GPS firmware & UART 9600 baud & 3.3 V logic expected & Each valid GGA sentence & RX GPIO 4, TX GPIO 5 in firmware \\
        \hline
        SHTC3 sensor bus & SHTC3 & ESP32 motor receiver & I2C 100 kHz & 3.3 V & 2 s & SDA GPIO 10, SCL GPIO 8 in prototype \\
        \hline
        Telemetry link & ESP32 GPS firmware & Flask ground station & Wi-Fi UDP JSON & RF link & Per GPS fix & Broadcast port 4210 \\
        \hline
        Cloud uplink & ESP32 GPS firmware & Hosted Flask server & HTTPS POST JSON & Wi-Fi & Per GPS fix when enabled & Uses \verb|X-API-Key| header \\
        \hline
        Camera stream & ESP32-CAM, phone, or webcam & Ground station / YOLO process & HTTP MJPEG & Network & Stream dependent & Registered through \verb|/api/camera| \\
        \hline
        Bench motor control & ESP32 receiver & H-bridge / motor driver & GPIO direction plus LEDC PWM & Driver dependent & 50 ms command loop & Stand-in for final ESC integration \\
        \hline
    \end{tabular}
    \caption{Electrical and software interface definitions.}
    \label{tab:interface_template}
\end{table}

\subsection{Technology (Languages, Hardware, etc.)}

FireFly combines embedded C/C++ firmware, Python backend services, computer vision, browser-based mapping, and 3D-printed mechanical design. The current embedded work targets ESP32-class microcontrollers using ESP-IDF. Implemented firmware includes UART GPS parsing, Wi-Fi station connection, UDP broadcast telemetry, optional HTTPS cloud upload, ESP-NOW command exchange, SHTC3 I2C temperature/humidity readout, and LEDC PWM motor-control bench tests.

\vspace{0.5em}

The ground station is written in Python using Flask. It stores live drone state in memory for local demonstrations and can optionally use SQLAlchemy/PostgreSQL with Supabase-style JWT authentication for registered drones and API keys. The browser UI uses Leaflet, JavaScript, and OpenStreetMap tiles to show drone markers, trails, bearing lines, fire markers, camera stream status, and active/stale alerts. The fire detection stack uses OpenCV and Ultralytics YOLOv8 weights trained for smoke and fire classes. The repository also contains a one-command launcher that starts the server, optional simulator, camera registration, and YOLO detector for local demonstrations.

\vspace{0.5em}

The long-term communications plan includes ESP-MESH or another MANET layer for self-organizing relay behavior, plus separate low-bandwidth and high-bandwidth links for telemetry and video. These links are design targets; the implemented prototype currently relies on local Wi-Fi UDP, HTTP APIs, ESP-NOW bench communication, and simulated drone motion.

\subsection{Simulations}

Before building and flying, we used simulation tools to verify that the FireFly configuration can sustain flight, carry its sensor payload, and meet the core mission requirement of delivering real-time environmental data during the pre-fire window. An online multicopter calculator (eCalc xcopterCalc) was used to evaluate combinations of frame size, motors, propellers, and battery against thrust-to-weight ratio, hover throttle, current draw, and expected flight time under sensor payload assumptions. This confirmed that the selected hardware provides sufficient thrust margin to carry the temperature/humidity sensor and smoke detector while maintaining acceptable hover efficiency and a flight time long enough to collect and transmit meaningful readings to the ground station.

\vspace{0.5em}

A simplified quadcopter dynamics model was used to tune attitude and altitude control parameters, verifying that the platform can hold a stable position over a monitoring area, a prerequisite for consistent sensor sampling. Research on quadcopter modeling confirms that such simulations can predict controller behavior and reduce the number of risky test flights needed to achieve stable flight. Together, these simulations provide evidence that FireFly can meet its need statement requirements: reaching areas with zero sensor infrastructure and sustaining flight long enough to collect and transmit temperature, humidity, and smoke data before committing to full implementation and field testing.


\subsection{Navigation and Perimeter Mapping}
FireFly uses an aspirational patrol-and-respond workflow designed to maximize aerial coverage, maintain communications, and focus sensing resources when a fire is detected. During the initial search phase, drones follow planned coverage paths such as a circular orbit or lawnmower pattern at a constant altitude. During this phase, the mature system will sample environmental data and stream imagery for fire detection. The current repository does not implement autonomous flight control; it simulates drone movement and validates the ground-station geolocation logic that would consume flight data.

\vspace{0.5em}

When the system detects a fire through visual inference or future environmental thresholds, it records telemetry containing GPS coordinates, altitude, heading, detection status, satellite count, and HDOP. The current ground station can then add fire markers directly from YOLO detections or compute a fire estimate by intersecting two recent drone bearing rays. This supports the eventual mapping sequence while keeping the present prototype focused on detection, telemetry, and geolocation rather than autonomous vehicle control.

\vspace{0.5em}

\subsubsection{Autonomy and Behavior}
The final flight logic is planned as a state machine that transitions between search, engagement, and recovery. During the search phase, the system will use a Boustrophedon, or lawnmower, path-planning pattern to ensure systematic coverage of a designated area with controlled sensor overlap. If a fire is detected, drones should maintain a safe standoff distance based on thermal intensity and operator-defined safety margins. In recovery, failsafes should handle return-to-home, signal loss, low battery, and hardware anomalies. The repository does not yet implement these flight-control states; it implements the telemetry and ground-station data products that the state machine would consume.

\vspace{0.5em}

\textbf{Search State:} In autonomous surveillance, the search state is the foundational operating mode where the system systematically monitors the geographic area to confirm a baseline and search for anomalies. The lawnmower pattern or horizontal sweep pattern is selected because it creates predictable coverage, simpler georeferencing, and repeatable sensor overlap.

The search state is characterized by steady-state flight at a constant altitude and velocity. This stability supports server-side processing by providing consistent video and reducing motion blur. Track spacing can be calculated from altitude and camera field of view, with overlap used to compensate for wind drift, airframe tilt, and inference delay. The current simulator uses circular motion rather than a lawnmower planner, so this section defines intended behavior for the next flight-control iteration.


Because the ground station receives video and telemetry continuously, predictable flight paths simplify georeferencing. When the server detects a flame, it can associate the detection with a known vehicle pose more reliably than it could during erratic movement. If a drone loses signal, the last known path also helps the operator estimate where to search.

\vspace{0.5em}

\textbf{Engagement and Mapping State:} This future phase shifts drones from broad surveillance to focused fire observation once a fire is confirmed. The primary observer continues reporting fire bearing and confidence, while secondary drones move to angular offsets around the fire to improve triangulation. The orbit or standoff path must protect the airframe from heat and updrafts and must include collision-avoidance spacing between drones.

\vspace{0.5em}

\textbf{Recovery State:} The recovery state is a high-priority override mode to protect the drone and people nearby. The intended behavior is immediate hover or return-to-home on telemetry loss, low battery, or hardware fault, followed by a controlled landing if recovery is not possible. Candidate thresholds are a soft battery failsafe near 30\% capacity and a critical landing threshold near 10\%, but these values must be validated with the final battery, ESC, and flight controller.

\subsubsection{Vision Detection}
The implemented vision pipeline runs on the ground-station computer rather than onboard the drone. The recorded-video detector loads YOLOv8 weights from \verb|Trained-Models/last.pt|, processes frames with OpenCV, labels detections as Smoke or Fire, and draws bounding boxes for operator review. A fire alert is only posted after the Fire class is present for at least 3 seconds, reducing false positives from brief frame-level detections. The posted fire payload includes confidence, estimated size, estimated area, and geolocation when drone state and camera geometry are available.

\vspace{0.5em}

The live-stream path uses MJPEG ingest. A camera URL can be registered with the ground station through \verb|/api/camera| using a camera-control token. The ingest worker reads frames from the MJPEG stream, runs the YOLO model every configured number of frames, and posts fire or smoke detections after a sustained 2 second window with a 5 second cooldown. If the selected drone has latitude, longitude, altitude, and heading, the detector estimates the fire bearing and ground intersection using a pinhole camera approximation; otherwise it reports the detection without coordinates.

\subsubsection{Multi-Drone Coordination}
The full FireFly system uses multiple drones for coverage, redundancy, and triangulation. During search, drones divide the survey region into lanes or orbits. When one drone detects a fire, it becomes the initial observer and continues reporting its position and bearing. Nearby drones are then assigned offset viewing positions so that at least two bearing rays intersect the same suspected ignition point. This geometry improves localization accuracy and helps reject false positives caused by a single camera angle.

\vspace{0.5em}

The current repository implements the coordination logic at the data-fusion level rather than in flight control. \verb|simulate_fleet.py| broadcasts two drone states over UDP. Each simulated drone marks \verb|fire\_detected| true when its heading falls within a camera field-of-view window around a fixed fire coordinate. The ground station waits for at least two fresh drones with fire detections and headings, then passes their latitude, longitude, and heading values into the triangulation function.

\subsubsection{Perimeter Generation and Area Calculation}
The proof-of-concept currently estimates a fire point, not a full perimeter. The triangulation module treats the local area as a flat-earth coordinate frame around a reference latitude and longitude. Each drone observation is converted into a ray using the drone position and heading. If two rays are not parallel and the intersection lies in front of the observing drone, the ground station converts the intersection back to latitude and longitude and creates a fire marker.

\vspace{0.5em}

The full perimeter-mapping design extends this point-estimation approach. As drones orbit or take offset views around the fire, the ground station will collect repeated boundary observations. Those observations can be ordered by bearing around the fire center, connected into a polygon, and converted to area using a local planar polygon-area calculation. YOLO bounding boxes and ground-sampling-distance estimates provide a first approximation of fire size in the current detector, but accurate perimeter mapping still requires calibrated camera geometry, synchronized timestamps, drone pose, and flight-test validation.

\subsubsection{Telemetry and Data}
The current telemetry packet is JSON and includes \verb|id|, \verb|seq|, \verb|lat|, \verb|lon|, \verb|alt_m|, \verb|heading_deg|, \verb|fire_detected|, \verb|sats|, \verb|hdop|, and \verb|ts_ms|. UDP broadcast uses port 4210 for local operation. Optional cloud upload posts the same JSON payload to the configured \verb|/api/drone| endpoint with an \verb|X-API-Key| header. The ground station exposes \verb|/api/state| for the dashboard, \verb|/api/fire| for fire detections, \verb|/api/camera| for camera stream registration, \verb|/api/triangulate| for manual triangulation, and \verb|/api/reset| for clearing local state.

\vspace{0.5em}

The future data-priority policy separates critical telemetry from high-bandwidth video. GPS, heading, battery, link health, and fire alerts should be transmitted first because they are required for safety and geolocation. Video frames, raw sensor logs, and dense perimeter observations can be throttled or dropped under poor link conditions. The current repository demonstrates this principle by allowing fire coordinates and confidence values to be posted independently of the displayed camera stream.



\newpage

% ============================================================
\section{Evaluation}
% ============================================================

\subsection{Functional Prototype}
The current FireFly prototype is strongest as a software and telemetry proof-of-concept. The repo can run a local ground station that displays live drone positions, trails, fire markers, bearing rays, and an optional camera panel. It accepts ESP32 GPS telemetry over UDP, supports optional authenticated cloud telemetry through \verb|/api/drone|, runs a simulated two-drone fleet, performs bearing-ray triangulation, and accepts YOLOv8 fire detections from recorded video or a live MJPEG stream.

\vspace{0.5em}

The hardware prototype has validated important bench-level subsystems but is not yet a fully integrated autonomous aircraft. ESP-NOW sender/receiver code demonstrates wireless command transfer, the motor receiver drives a bench motor through GPIO direction pins and LEDC PWM, and SHTC3 temperature/humidity values are read over I2C and returned to the sender. A DC brushed motor and H-bridge were used as stand-ins while waiting for final ESC and brushless motor integration. The CAD frame, mass estimates, BOM, and circuit diagrams define the intended airframe and electronics direction.

\subsection{Testing}
Testing is divided into software validation, bench validation, and future flight validation. Software validation covers the ground station, simulated fleet, YOLO detector, MJPEG ingest, and fire posting APIs. Bench validation covers ESP-NOW communication, SHTC3 readout, potentiometer command input, and PWM motor response. Future validation must measure flight endurance, outdoor telemetry range, calibrated environmental thresholds, camera geolocation accuracy, and multi-drone field behavior.

\subsection{Measured Results Summary}

\begin{table}[H]
    \centering
    \small
    \begin{tabular}{|p{0.24\textwidth}|p{0.18\textwidth}|p{0.18\textwidth}|p{0.24\textwidth}|}
        \hline
        \textbf{Metric} & \textbf{Measured Value} & \textbf{Units} & \textbf{Test Context / Notes} \\
        \hline
        Hover time & Not yet measured & min & Flight stack not integrated; eCalc/design estimate is 5--8 min for Option A \\
        \hline
        Telemetry range & Not yet measured outdoors & m & Current proof-of-concept uses local Wi-Fi UDP and simulator \\
        \hline
        Packet latency & Not yet characterized & ms & Dashboard polls \verb|/api/state| every 500 ms; UDP packet timing still needs logging \\
        \hline
        Sensor update period & 2 & s & SHTC3 bench firmware sends temperature/humidity every 2 seconds \\
        \hline
        Alert trigger delay & 3 / 2 & s & Recorded-video detector uses 3 s sustained fire; MJPEG ingest uses 2 s sustained fire/smoke \\
        \hline
        Total all-up weight & 190--230 & g & Estimated Option A baseline; scale measurement still required \\
        \hline
    \end{tabular}
    \caption{Measured and estimated prototype results.}
    \label{tab:measured_results_template}
\end{table}

\subsection{Observed Issues and Next Iteration}

\begin{table}[H]
    \centering
    \small
    \begin{tabular}{|p{0.22\textwidth}|p{0.20\textwidth}|p{0.20\textwidth}|p{0.18\textwidth}|}
        \hline
        \textbf{Issue} & \textbf{Observed During} & \textbf{Probable Cause} & \textbf{Planned Design Change} \\
        \hline
        Motor pod sizing mismatch & Physical prototype and received motor fit check & CAD tolerance and motor dimension assumptions & Update motor-pod dimensions and reprint fit-check section \\
        \hline
        Brushless motor integration blocked & Propulsion bench phase & ESC delivery/selection delay & Continue H-bridge bench tests, then integrate final 4-in-1 ESC \\
        \hline
        ESP-NOW motor command path lacks authentication and clamping in active receiver & Source review & Prototype code prioritizes quick communication demo & Validate peer MAC, clamp duty values, and add authenticated/encrypted command path before flight \\
        \hline
        Autonomous flight not implemented & Repo review & Project scope shifted toward geolocation pipeline & Treat flight automation as future work and validate geolocation independently \\
        \hline
    \end{tabular}
    \caption{Observed issues and next iteration actions.}
    \label{tab:issue_template}
\end{table}

\newpage

% ============================================================
\section{Parts List}
% ============================================================

\begin{table}[H]
    \centering
    \small
    \begin{tabular}{|l|c|c|}
        \hline
        \textbf{Part} & \textbf{Supplier} & \textbf{Quantity} \\
        \hline
        CarbonX PLA+CF Filament & Creality & 1 \\
        \hline
        Battery: LiPo 850 mAh 3S 25/30C & HobbyKing & 1 \\
        \hline
        Motor: SunnySky R1106 FPV Brushless Motors & Druav & 4 \\
        \hline
        Propeller: Gemfan Hurricane 2009-3 Tri-Blade (Set of 4) & Pyrodrone & 1 \\
        \hline
        HAKRC 20A 4-in-1 \verb|3-6S BLHeli_32| 5V 3A BEC DShot1200 Mini ESC & Alibaba & 1 \\
        \hline
        Perforated Circuit Board & Adafruit & 2 \\
        \hline
        Female Header Pins 2x6 & AliExpress & 4 \\
        \hline
        Female Header Pins 2x8 & AliExpress & 2 \\
        \hline
        Female Header Pins 2x10 & Mikroe & 2 \\
        \hline
    \end{tabular}
    \caption{Itemized Parts List}
    \label{tab:budget}
\end{table}

\subsection{Expanded Bill of Materials}

\begin{table}[H]
    \centering
    \scriptsize
    \begin{tabular}{|p{0.20\textwidth}|p{0.16\textwidth}|p{0.10\textwidth}|p{0.10\textwidth}|p{0.10\textwidth}|p{0.12\textwidth}|p{0.12\textwidth}|}
        \hline
        \textbf{Part / Part Number} & \textbf{Function} & \textbf{Qty.} & \textbf{Unit Cost} & \textbf{Total} & \textbf{Source / Lead Time} & \textbf{Status} \\
        \hline
        LiPo 850 mAh 3S 25/30C & Option A battery & 1 & \$6.18 & \$6.18 & HobbyKing & Selected \\
        \hline
        SunnySky R1106 FPV Brushless Motor & Option A propulsion & 4 & \$11.99 & \$47.96 & Druav & Selected \\
        \hline
        Gemfan Hurricane 2009-3 Tri-Blade & Option A propellers & 1 set & \$4.99 & \$4.99 & Pyrodrone & Selected \\
        \hline
        HAKRC 20A 4-in-1 BLHeli\_32 ESC & Motor control / BEC & 1 & \$30.12 & \$30.12 & Alibaba & Delayed / integration pending \\
        \hline
        Perforated circuit board & Electronics mounting & 2 & \$4.95 & \$9.90 & Adafruit / BELS & Received \\
        \hline
        Female header pins 2x6 & Board interconnect & 4 & \$0.17 & \$0.68 & eBay / BELS & Received \\
        \hline
        Female header pins 2x8 & Board interconnect & 2 & \$0.09 & \$0.18 & AliExpress / BELS & Received \\
        \hline
        Female header pins 2x10 & Board interconnect & 2 & \$0.72 & \$1.44 & Mikroe / BELS & Received \\
        \hline
        CADDXFPV Ant Lite Analog Camera & Optional video payload & 1 & \$15.90 & \$15.90 & CaddxFPV & Optional / compatibility risk \\
        \hline
        TFMini-S Micro LiDAR Module & Optional range payload & 1 & \$72.95 & \$72.95 & SparkFun & Optional future sensor \\
        \hline
        3S 1000 mAh battery & Option B alternate battery & 1 & \$10.99 & \$10.99 & HobbyKing & Alternate \\
        \hline
        RCINPOWER GTS 1204 5000KV motor & Option B alternate motor & 4 & \$18.99 & \$75.96 & Pyrodrone & Alternate \\
        \hline
        HQProp T3.5x2x3 propeller & Option B alternate propellers & 1 set & \$3.99 & \$3.99 & GetFPV & Alternate \\
        \hline
    \end{tabular}
    \caption{Expanded bill of materials from the latest repository BOM.}
    \label{tab:bom_expanded_template}
\end{table}

\subsection{Cost Summary}

\begin{table}[H]
    \centering
    \small
    \begin{tabular}{|p{0.28\textwidth}|p{0.18\textwidth}|p{0.34\textwidth}|}
        \hline
        \textbf{Category} & \textbf{Subtotal} & \textbf{Notes} \\
        \hline
        Mechanical / printed parts & Not priced & CarbonX PLA+CF filament is listed but no unit cost is recorded in the repo BOM \\
        \hline
        Propulsion and power & \$89.25 & Option A battery, motors, propellers, and ESC \\
        \hline
        Computing and telemetry & \$12.20 & Perforated boards and headers; ESP32 boards appear to be reused/lab inventory \\
        \hline
        Sensors and payload & \$88.85 & Optional camera and LiDAR payloads \\
        \hline
        Spares / iteration cost & \$90.94 & Option B alternate propulsion set if selected \\
        \hline
        Total prototype cost & \$190.30 & Option A plus listed optional camera/LiDAR and electronics hardware, excluding filament and reused ESP32 boards \\
        \hline
    \end{tabular}
    \caption{Prototype cost summary.}
    \label{tab:cost_summary_template}
\end{table}

\subsection{Procurement and Assembly Notes}
\begin{itemize}
    \item Long-lead items and alternates: the ESC was a schedule blocker; Option B propulsion parts are documented as alternates.
    \item Components reused from lab inventory: perforated boards, female headers, H-bridge bench hardware, and ESP32 development boards were sourced from BELS or existing lab resources.
    \item Components still unselected or unverified: final smoke/gas sensor, calibrated fire-risk thresholds, battery monitor, production camera path, and mesh/long-range radio hardware.
    \item Assembly dependencies that affect schedule: motor pod fit, ESC integration, battery mounting, wire routing, sensor mounting, and safe tethered-flight setup.
\end{itemize}

\newpage

% ============================================================
\section{References}
% ============================================================

\begin{enumerate}
    \item Espressif Systems, ``ESP-NOW,'' ESP-IDF Programming Guide. \url{https://docs.espressif.com/projects/esp-idf/en/latest/esp32c6/api-reference/network/esp_now.html}
    \item Espressif Systems, ``ESP-WIFI-MESH,'' ESP-IDF Programming Guide. \url{https://docs.espressif.com/projects/esp-idf/en/release-v4.4/esp32/api-guides/esp-wifi-mesh.html}
    \item NASA Earthdata FIRMS, ``Active Fire Data Attributes for MODIS and VIIRS.'' \url{https://www.earthdata.nasa.gov/data/tools/firms/active-fire-data-attributes-modis-viirs}
    \item NASA Earthdata, ``VIIRS I-Band 375 m Active Fire Data.'' \url{https://www.earthdata.nasa.gov/data/instruments/viirs/viirs-i-band-375-m-active-fire-data}
    \item ALERTCalifornia, ``About.'' \url{https://alertcalifornia.org/about/}
    \item Federal Aviation Administration, ``Small Unmanned Aircraft Systems (UAS) Regulations (Part 107).'' \url{https://www.faa.gov/newsroom/small-unmanned-aircraft-systems-uas-regulations-part-107}
    \item Ultralytics, ``Ultralytics YOLO Docs.'' \url{https://docs.ultralytics.com/}
    \item Thilak R., ``Forest-fire-detection-using-YOLOv8.'' \url{https://github.com/thilak-r/Forest-fire-detection-using-YOLOv8}
    \item eCalc, ``xcopterCalc.'' \url{https://ecalc.ch/xcoptercalc.php}
\end{enumerate}

% ============================================================
\section{Appendix 1 - Problem Formulation}
% ============================================================

\subsection{Conceptualisations}

FireFly is a lightweight multi-drone wildfire monitoring concept designed to collect and transmit useful fire intelligence in remote areas with no existing infrastructure. The system is built around four functions: aerial observation, environmental sensing, resilient telemetry, and map-based decision support.

\vspace{0.5em}

The intended drone carries an onboard microcontroller, GPS, camera, environmental sensors, flight controller, propulsion system, and LiPo battery. In the current repository, GPS telemetry, temperature/humidity readout, ESP-NOW communication, motor bench control, camera ingest, and ground-station mapping are implemented as separate proof-of-concept subsystems rather than a fully integrated flight vehicle.

\vspace{0.5em}

Data is transmitted to a Python Flask ground station. The current local path uses Wi-Fi UDP broadcast on port 4210, while the optional cloud path posts authenticated JSON telemetry to \verb|/api/drone|. The browser dashboard displays drone state, fire markers, camera stream status, and threshold or detection alerts.

\subsection{Brainstorming}
Several frame geometries were evaluated including X and H-style quadcopter
layouts. An X-frame was selected for its balanced motor distribution,
stability, and ease of sensor mounting.

\vspace{0.5em}

The sensing approach was inspired by existing ground-based forest sensor
networks, which are geographically limited and difficult to scale into remote
terrain. FireFly was conceived as a mobile alternative: a drone that could
collect the same environmental data in areas where fixed infrastructure does
not exist.

\vspace{0.5em}

For communication, Bluetooth Low Energy, Wi-Fi, ESP-NOW, and mesh-style networking were evaluated. ESP-NOW was selected for early microcontroller-to-microcontroller bench testing because it is connectionless and low overhead. Wi-Fi UDP and HTTPS were selected for the current GPS-to-ground-station proof-of-concept because they integrate directly with the Flask dashboard. ESP-MESH, LoRa, and COFDM remain future candidates for field-ready range and relay behavior.

\vspace{0.5em}

Power and propulsion options were compared using eCalc to evaluate
thrust-to-weight ratio, hover current, and projected flight time across
different motor, propeller, and LiPo configurations. This guided selection
of a 3S LiPo and motor pairing that balances flight endurance with the
electrical load of the sensor and communication hardware.

\subsection{Decision Tables}

The following decision tables evaluate key design alternatives across weighted criteria. Scores are rated 1--5 per criterion, multiplied by the criterion weight, and summed to produce a weighted total.

\begin{table}[H]
\centering
\small
\begin{tabular}{|l|c|c|c|c|}
\hline
\textbf{Criterion} & \textbf{Weight} & \textbf{X-Frame} & \textbf{H-Frame} & \textbf{+ Frame} \\
\hline
Stability & 25\% & 5 & 3 & 4 \\
\hline
Sensor Mounting & 20\% & 5 & 4 & 3 \\
\hline
Ease of Assembly & 20\% & 4 & 4 & 3 \\
\hline
Printability & 20\% & 4 & 3 & 5 \\
\hline
Repairability & 15\% & 4 & 5 & 3 \\
\hline
\textbf{Weighted Total} & & \textbf{4.45} & \textbf{3.75} & \textbf{3.65} \\
\hline
\end{tabular}
\caption{Frame Geometry Decision Table}
\end{table}

\begin{table}[H]
\centering
\small
\begin{tabular}{|l|c|c|c|c|}
\hline
\textbf{Criterion} & \textbf{Weight} & \textbf{Peer-to-Peer Wireless} & \textbf{Wi-Fi} & \textbf{BLE} \\
\hline
Latency & 30\% & 5 & 3 & 4 \\
\hline
Range & 25\% & 5 & 4 & 2 \\
\hline
Infrastructure Required & 25\% & 5 & 2 & 4 \\
\hline
Implementation Complexity & 20\% & 4 & 3 & 4 \\
\hline
\textbf{Weighted Total} & & \textbf{4.80} & \textbf{3.00} & \textbf{3.50} \\
\hline
\end{tabular}
\caption{Communication Protocol Decision Table}
\end{table}

\begin{table}[H]
\centering
\small
\begin{tabular}{|l|c|c|c|}
\hline
\textbf{Criterion} & \textbf{Weight} & \textbf{PLA} & \textbf{Carbon Fiber PLA} \\
\hline
Weight & 30\% & 4 & 5 \\
\hline
Printability & 25\% & 5 & 3 \\
\hline
Structural Rigidity & 25\% & 3 & 5 \\
\hline
Repairability & 20\% & 5 & 3 \\
\hline
\textbf{Weighted Total} & & \textbf{4.20} & \textbf{4.00} \\
\hline
\end{tabular}
\caption{Frame Material Decision Table}
\end{table}

\subsection{Morphological Charts}
The following combinations were explored for the primary subsystems:
\begin{figure}[H]
    \centering
    \includegraphics[width=0.5\linewidth, height=0.45\textheight]{images/morphological.png}
    \caption{Morphological analysis chart outlining potential components of FireFly}
    \label{fig:morphological}
\end{figure}

\newpage

% ============================================================
\section{Appendix 2 - Planning}
% ============================================================

\subsection{Basic Plan / Gantt Chart}
\begin{figure}[H]
    \centering
    \includegraphics[width=1\linewidth, height=0.50\textheight, keepaspectratio]{images/Gantt_Chart_v2.png}
    \caption{Gantt Chart}
    \label{fig:gantt_chart}
\end{figure}

\subsection{Division of Labor}

\begin{table}[H]
    \centering
    \begin{tabular}{|l|l|}
        \hline
        \textbf{Team Member} & \textbf{Primary Responsibility} \\
        \hline
        Christian Li & Flight Controller Programming \\
        \hline
        Kenny Wang & Hardware Integration \\
        \hline
        Kevin Yao & Mechanical Design / CAD \\
        \hline
        Matin Sobhani & Mechanical Design / CAD \\
        \hline
        Nick Lafredo & Power Systems / Documentation \\
        \hline
        Parsa Rezvani & Testing and Validation / Telemetry \\
        \hline
        Ryan Yu & Flight Controller Programming \\
        \hline
    \end{tabular}
    \caption{Division of Labor during Prototyping}
\end{table}

\subsection{Collaboration}
\begin{itemize}
    \item \textbf{Github:} CI/CD Tool
    \item \textbf{Discord:} Communication and meetings
    \item \textbf{Notion:} Gantt Chart
    \item \textbf{Trello:} Kanban board for task tracking
\end{itemize}

\newpage

% ============================================================
\section{Appendix 3 - Test Plan \& Results}
% ============================================================

The test plan consists of three phases:
\begin{enumerate}
    \item \textbf{Bench Testing:}

    We began by creating a prototype circuitry system that utilizes an H-Bridge in substitution for the Electronic Speed Controller, along with a DC brushed motor in place of the designed DC brushless motor. This system was used to test the wireless communication between the ground-side and drone-side microcontrollers, allowing us to develop and validate the data transmission pipeline while awaiting final parts.

    \item \textbf{Tethered Flight:}

    We will integrate the ESC and DC brushless motors into the physical 3D-printed frame and test flight capability in a controlled indoor environment. During this phase we will emulate conditions that would flag a fire danger zone (elevated temperature, low humidity, and elevated smoke readings) and verify that threshold alerts are correctly triggered and transmitted to the ground station.

    \item \textbf{Open Field Flight:}

    We will test the drone in an open field to evaluate durability, structural integrity, and real-world sensor performance. Outdoor conditions including wind and heat will allow us to assess how environmental factors affect flight stability and sensor accuracy, providing direct feedback on both the mechanical design and the telemetry pipeline under realistic deployment conditions.
\end{enumerate}

\subsection{Bench Test Log}

\begin{table}[H]
    \centering
    \small
    \begin{tabular}{|p{0.12\textwidth}|p{0.18\textwidth}|p{0.18\textwidth}|p{0.18\textwidth}|p{0.12\textwidth}|p{0.12\textwidth}|}
        \hline
        \textbf{Test ID} & \textbf{Objective} & \textbf{Procedure} & \textbf{Expected Result} & \textbf{Observed Result} & \textbf{Pass / Fail} \\
        \hline
        B-01 & Verify ESP-NOW command path & Send potentiometer command from remote ESP32 to receiver ESP32 & Receiver updates motor direction and PWM duty & Bench prototype demonstrated direction and speed control with H-bridge/DC motor & Pass for bench prototype \\
        \hline
        B-02 & Verify temperature/humidity readout & Read SHTC3 over I2C on receiver and return values to sender & Serial output shows temperature and relative humidity & Firmware implements CRC-checked SHTC3 read and ESP-NOW return packet & Pass for readout; calibration pending \\
        \hline
        B-03 & Verify GPS telemetry packet shape & Parse GPS GGA sentence and broadcast JSON to UDP port 4210 & Ground station receives id, seq, lat, lon, altitude, heading, sats, and HDOP & Firmware and Flask UDP listener implement the packet path & Implemented; range pending \\
        \hline
        B-04 & Verify visual detection to map path & Run YOLO detector on fire video or MJPEG stream and post to server & Sustained fire detection creates map marker & Detector and \verb|/api/fire| route implement this behavior & Software proof complete \\
        \hline
    \end{tabular}
    \caption{Bench-test log and proof-of-concept status.}
    \label{tab:bench_test_template}
\end{table}

\subsection{Tethered Flight Plan}

\begin{table}[H]
    \centering
    \small
    \begin{tabular}{|p{0.12\textwidth}|p{0.14\textwidth}|p{0.12\textwidth}|p{0.14\textwidth}|p{0.14\textwidth}|p{0.18\textwidth}|}
        \hline
        \textbf{Flight ID} & \textbf{Configuration} & \textbf{Duration} & \textbf{Stability Notes} & \textbf{Telemetry Notes} & \textbf{Issues / Corrective Action} \\
        \hline
        T-01 & Frame, final ESC, brushless motors, no sensor payload & Short hover bursts & Verify arm rigidity, motor direction, and throttle response & Confirm ground station continues to receive telemetry if GPS is mounted & Planned after ESC integration \\
        \hline
        T-02 & Frame plus ESP32 telemetry payload & 1--2 min tethered hover & Check vibration impact on wiring and GPS heading stability & Log packet sequence continuity and stale-drone behavior & Planned \\
        \hline
        T-03 & Frame plus camera or simulated camera payload & 1--2 min tethered hover & Verify payload mounting and center of mass & Verify camera URL registration and map UI responsiveness & Planned \\
        \hline
    \end{tabular}
    \caption{Tethered-flight validation plan.}
    \label{tab:tethered_flight_template}
\end{table}

\subsection{Open Field Flight Plan}

\begin{table}[H]
    \centering
    \small
    \begin{tabular}{|p{0.10\textwidth}|p{0.12\textwidth}|p{0.12\textwidth}|p{0.12\textwidth}|p{0.12\textwidth}|p{0.14\textwidth}|p{0.18\textwidth}|}
        \hline
        \textbf{Flight ID} & \textbf{Wind / Weather} & \textbf{Payload} & \textbf{Flight Time} & \textbf{Range} & \textbf{Alert Behavior} & \textbf{Observations} \\
        \hline
        F-01 & Calm, clear & GPS telemetry only & Target 3--5 min & Local Wi-Fi range walk & No fire source; telemetry only & Establish safe baseline flight \\
        \hline
        F-02 & Calm, clear & GPS plus camera stream & Target 3--5 min & Measure stream reliability & Use controlled visual target or recorded stream & Validate operator map workflow \\
        \hline
        F-03 & Light wind limit & GPS plus environmental payload & Target 5 min & Measure packet loss with movement & Inject safe heat/smoke substitute if permitted & Validate threshold workflow without open flame \\
        \hline
    \end{tabular}
    \caption{Open-field flight validation plan.}
    \label{tab:open_field_template}
\end{table}

\subsection{Sensor Calibration Plan}

\begin{table}[H]
    \centering
    \small
    \begin{tabular}{|p{0.16\textwidth}|p{0.16\textwidth}|p{0.16\textwidth}|p{0.16\textwidth}|p{0.18\textwidth}|}
        \hline
        \textbf{Sensor} & \textbf{Reference Instrument} & \textbf{Reference Value} & \textbf{Measured Value} & \textbf{Offset / Correction} \\
        \hline
        SHTC3 temperature & Calibrated thermometer & Ambient lab temperature & ESP32 serial reading & Compute temperature offset before field use \\
        \hline
        SHTC3 humidity & Calibrated hygrometer & Ambient relative humidity & ESP32 serial reading & Compute humidity offset before field use \\
        \hline
        Camera geolocation & Measured ground marker & Known marker position & YOLO/bearing estimate & Tune camera HFOV, VFOV, tilt, and altitude assumptions \\
        \hline
    \end{tabular}
    \caption{Sensor-calibration plan.}
    \label{tab:calibration_template}
\end{table}

\subsection{Post-Test Design Changes}
\begin{itemize}
    \item Motor fit issue: update CAD motor-pod dimensions and reprint before brushless propulsion testing.
    \item ESC delay: keep the H-bridge/DC motor rig for command-path testing, then repeat with the selected 4-in-1 ESC.
    \item ESP-NOW command risk: validate source MAC, clamp duty values, and add authentication before any powered flight test.
    \item Geolocation scope: continue validating map, camera, and triangulation independently while flight automation remains future work.
\end{itemize}

\newpage

% ============================================================
\section{Appendix 4 - Review}
% ============================================================

\textit{Reflection of each team member's thoughts on the project.}

\vspace{1em}

\textbf{Christian Li:} Looking back on this project, one of the biggest lessons I learned is the importance of starting with a clearly defined need rather than a preconceived solution. We approached the project wanting to build a drone, and only later did we step back and ask why a drone was the right solution for the problem. If we had focused earlier on the lack of fire risk data in remote areas, we could have identified the required characteristics sooner and avoided some of the design and task adjustments we made midway through the quarter. Even so, our team communicated well and adapted quickly. The flexibility helped us stay on track and produce a coherent design. This experience also gave me valuable insight into working as a team, improving how we communicate, distribute work, and learn new tools and environments together.

\vspace{1em}

\textbf{Kenny Wang:} Restarting this project I think that we should focus on not building a product that interests us but develop the idea based on a universal problem or a need. From there we can then focus on implementing it in a way that can handle this universal problem or need. For example, if we had started with the problem being that we have no way of reaching remote areas with no fire detection, we can then build upon that and not just specifically building a flying drone but maybe a rover, or even something that does not even have to move. There would be a variety of implementations of helping alleviate the same problem and need while having numerous amounts of angles to implement that solution. Even with some setbacks that we faced during the beginning stages of the project we were able to have good communication within the team to implement these changes so that we were able to complete our design work.

\vspace{1em}

\textbf{Kevin Yao:} Reflecting back to the beginning of this project, I think we should have focused on the process instead of the end goal of the design we wanted to make. One thing that we did well was forming our group and communicating often. We were vocal on what we wanted to do and had ideas on what the design would look like. A problem that occurred was our need and goal statement. What we should have done differently was to define clear objectives of what needed to be done and check the requirements. Another thing we should have done earlier was setting up the github with our other tools to have a cleaner pipeline of our progress. Although we had some setbacks, our team was able to implement these changes quickly and improved our communication and workflow.

\vspace{1em}

\textbf{Matin Sobhani:} If we were starting this project again, I would make sure we had a solid need/goal statement, and verified it with the TA and professor. We thought we had a solid one early on around week 4, but ended up having to change it around week 8. Thankfully this did not affect us too negatively since our underlying prototype idea stayed mostly the same, but it did cost us time and forced us to reframe parts of our documentation. In hindsight, we should have scheduled a formal requirement check-in earlier and created clearer milestones tied to those requirements so we could catch misalignment sooner. What went well is that we stayed very communicative throughout the project, kept each other up-to-date, and adapted quickly whenever priorities changed.

\vspace{1em}

\textbf{Nick Lafredo:} This project is the first I've done on such a large scale, and it taught me a lot about the timeline and organization needed to accomplish a group goal. If we were to go back and start a new project, we would focus more on the need and solution, rather than starting with a product we wanted to make and finding a need that fits it. It is also now much clearer that we are not designing the prototype, but creating a general design that could be implemented in the future, with the prototype proving its possibility. We faced a few setbacks with our design and deciding what problem we were aiming to solve, but ultimately, it made our final idea much stronger and well-defined. Going forward, we will have more frequent check ins with our team and our instructors to make sure we stay on the right track. We now have a clear path towards completing the design, and I am confident that we will reach our goal.

\vspace{1em}

\textbf{Parsa Rezvani:} This project has been a significant learning experience for both me and the team. Initially, we approached it as an exercise in selecting an interesting idea and producing documentation around it. Over time, however, we came to understand that the class demanded something closer to real-world engineering practice: defining a defensible need, building toward a clear goal, and maintaining a concept that could stand on its own merits in a professional context. Had we internalized that framing earlier, our process would have been more focused from the start. Despite the adjustments along the way, I am proud of how the team collaborated and adapted. With our components arriving soon, I am confident we are well positioned to make meaningful progress on the prototype.

\vspace{1em}

\textbf{Ryan Yu:} Looking back at the beginning of this project, there are a few things that went well and a few things that went poorly. I think we got our group together, talked about initial parts, calculations, and ideas quite early, and were on top of things. However, I think the communication between the professor, TA, and our group was extremely poor. It was not clearly laid out that we were designing a product but producing a prototype. We had difficulties having useful long discussions with the professor and were often confused on what he wanted, constantly changing last minute.

\end{document}
